\section{Selection Sort: Fragments, a Summary, and an Induction}\label{sec:sel}

\subsection{The program}

\begin{lstlisting}
method main(a: array<int, int>, n: int) returns (b: array<int, int>)
  requires 0 <= n
  ensures sorted(b, 0, n)
  ensures perm(b, a)
{
  var i: int; var j: int; var m: int; var t: int;
  b := a; i := 0;
  while i < n {
    m := i; j := i + 1;
    while j < n {
      if b[j] < b[m] { m := j; }
      j := j + 1;
    }
    t := b[i]; b[i] := b[m]; b[m] := t;
    i := i + 1;
  }
}
\end{lstlisting}

\noindent
Arrays are total maps $\mathsf{array}\langle\mathsf{int},\mathsf{int}\rangle$; $\sorted(a, lo, hi)$
states that $a[lo..hi)$ is nondecreasing, and $\perm(a, b)$ that $a$ is a permutation of $b$ as a total
map, i.e. $b$ composed with a bijection of the indices. Both are abstract predicates, declared without
a definition; their properties are the lemmas below. Whole-array permutation needs no side conditions
on the ranges of swapped positions: a swap at any two positions preserves it.

\subsection{The trace equation system}

Symbolic execution of the outer loop (Definition 4.4 of the paper) introduces the synthetic counter
$K = \cnt(0, 1, n')$ of $n'$ iterations and one trace per modified variable. The inner loop is
summarized during the execution of the body (Theorem 4.8): its own trace equation system is in terms of
fresh variables, which become fresh \emph{sequences} over the outer iterations, and it is recorded as the
iteration constraint of the outer loop. After \texttt{-simplify} distributes the constraint over its
conjuncts, the goal is (variables renamed, the fragments of the inner loop grouped):
\begin{align}
  K &= \cnt(0, 1, n') \qquad I = \cnt(0, 1, n') \qquad \intf(n') = n \qquad b = \cur(B) \label{eq:outer}\\
  B &= \tr{a}{\bl{K}, \bl{I}, \seq{M}}{x, k, i, m \mapsto \swap{x}{k}{m}} \label{eq:B}\\[2pt]
  &\al{\bl{K}, \bl{I}, \seq{M}, \seq{M}'}{k, i, m, M \mapsto m = \cur(M)} \label{eq:f1}\\
  &\al{\bl{K}, \bl{I}, \bl{B}, \seq{M}', \seq{J}'}{k, i, x, M, J \mapsto {} \nonumber\\
  &\qquad\qquad M = \tr{k}{\bl{J}}{p, j \mapsto \upd{x[j] < x[p]}{j}\ \ \upd{x[p] \le x[j]}{p}}} \label{eq:f2}\\
  &\al{\bl{K}, \bl{I}, \seq{N}, \seq{K}'}{k, i, n'', K' \mapsto K' = \cnt(0, 1, n'')} \label{eq:f3}\\
  &\al{\bl{K}, \bl{I}, \seq{K}', \seq{J}'}{k, i, K', J \mapsto J = \sums(k + 1, \mapc{\bl{K'}}{1})} \label{eq:f4}\\
  &\al{\bl{K}, \bl{I}, \seq{J}', \seq{J}}{k, i, J, j \mapsto j = \cur(J)} \label{eq:f5}\\
  &\al{\bl{K}, \bl{I}, \seq{J}'}{k, i, J \mapsto \al{\bl{J}}{j \mapsto j < n}} \label{eq:f6}\\
  &\al{\bl{K}, \bl{I}, \seq{J}}{k, i, j \mapsto n \le j} \label{eq:f7}\\[2pt]
  &\Longrightarrow\ \sorted(b, 0, n) \wedge \perm(b, a) \nonumber
\end{align}
Equation \eqref{eq:B} is the outer array trace: in iteration $k$ (with $i = k$, since $I$ and $K$ are
equal counters that the simplifier does not identify), the positions $k$ and $m$ are swapped, where $m$ is
the result of the inner loop, the entry of the sequence $\seq{M}$ at $k$. The fragments
\eqref{eq:f1}--\eqref{eq:f7} describe that inner loop: its result $m$ is the current value of a
\emph{selection trace} $M$ \eqref{eq:f1}, which starts at $k$ and moves to $j$ whenever $x[j] < x[p]$,
over the inner counter $J$ \eqref{eq:f2}; $J$ is the trace of $j$, increased by 1 from $k + 1$ over the
synthetic inner counter $K'$ of $n''$ steps \eqref{eq:f3}, \eqref{eq:f4}, with the final value $j$
\eqref{eq:f5}; the test $j < n$ holds in every inner iteration \eqref{eq:f6} and fails at the exit
\eqref{eq:f7}. The generic fold of \texttt{teq.dfy} already turned the trace of $j$ into
$\sums(k + 1, \mapc{\bl{K'}}{1})$, since inside a fragment the equation $K' = \cnt(0, 1, n'')$ of the
sibling fragment \eqref{eq:f3} is not followed (problem~1 of nested loops in \texttt{NOTES.md}).

\subsection{Step 1: the inner counter (generic fragment lemmas)}\label{sec:sel:frag}

The two lemmas of \texttt{fragments.dfy} relate fragments of one inner loop to each other. They are
rewrite rules whose left-hand side is a fragment and whose side condition is a sibling fragment, matched
against the facts of the goal, which binds the sequences that the left-hand side does not mention.
\begin{align*}
  &\lemma{invs\_sums\_counter}:\\
  &\quad \al{ks, Ks, Ns}{k, K, n'' \mapsto K = \cnt(0, 1, n'')}\ \Longrightarrow \\
  &\al{ks, Ks, Js}{k, K, J \mapsto J = \sums(X(k), \mapc{\bl{K}}{c})}\\
  &\quad\Longleftrightarrow\ \al{ks, Ns, Js}{k, n'', J \mapsto J = \cnt(X(k), c, n'')}\\[4pt]
  &\lemma{invs\_current\_counter}:\\
  &\quad \al{ks, Ns, Js}{k, n'', J \mapsto J = \cnt(X(k), c, n'')}\ \Longrightarrow \\
  &\al{ks, Js, js}{k, J, j \mapsto j = \cur(J)}\\
  &\quad\Longleftrightarrow\ \al{ks, Ns, js}{k, n'', j \mapsto j = X(k) + \intf(n'') \cdot c}
\end{align*}
The initial value $X(k)$ is a higher-order pattern variable (Miller fragment): it matches $0$, $k$,
$k + 1$, or the value of a further dependency, so one lemma covers all inner counters. The first lemma
matches fragment \eqref{eq:f4} with $ks := \bl{K}$, $Ks := \seq{K}'$, $Js := \seq{J}'$,
$X := \lambda k.\, k + 1$, $c := 1$; its side condition is found as fragment \eqref{eq:f3}, which binds
$Ns := \seq{N}$. The unused dependency $\bl{I}$ of the fragment is selected by no pattern dependency and
dropped. The result replaces \eqref{eq:f4}:
\begin{equation}
  \al{\bl{K}, \seq{N}, \seq{J}'}{k, n'', J \mapsto J = \cnt(k + 1, 1, n'')}. \label{eq:f4b}
\end{equation}
The second lemma matches \eqref{eq:f5} with $Js := \seq{J}'$, $js := \seq{J}$, finds \eqref{eq:f4b} as its
side condition, and yields
\begin{equation}
  \al{\bl{K}, \seq{N}, \seq{J}}{k, n'', j \mapsto j = \intf(n'') + k + 1}. \label{eq:f5b}
\end{equation}
Conceptually, the fragment lemmas do by rewriting what the simplifier would do if the fragments of one
inner loop shared their facts (the sharing was implemented and reverted, since it made rewriting cycle):
they propagate the definition of a bound variable from one fragment into another. They are candidates for
\texttt{teq.dfy}; they close the two minimal nested-loop problems \texttt{invs\_counter\_equation.dfy} and
\texttt{invs\_fragments.dfy}, and \texttt{nested\_sum.dfy}, \texttt{nested\_mult.dfy},
\texttt{distinct.dfy} and the sorting algorithms use them.

\subsection{Step 2: the summary of the inner loop}\label{sec:sel:sum}

The inner loop finds the position of the minimum of $x[k..n)$. We name that position by an abstract
function $\argmin(x, lo, hi)$, declared without a definition, and state the summary as a rewrite rule
whose side conditions are the fragments, in the forms they have after Step~1:
\begin{align*}
  &\lemma{current\_selection\_swaps}:\\
  &\quad \al{ks, ms, Ms}{k, m, M \mapsto m = \cur(M)} \\
  &\quad\wedge\ \mathsf{all}(\mathsf{for}\ ks, \bl{\tr{a}{ks, is, ms}{x, k, i, m \mapsto \swap{x}{k}{m}}}, Ms, Js\\
  &\qquad\qquad \mathsf{check}\ k, x, M, J \mapsto {}\\
  &\qquad\qquad\qquad M = \tr{k}{\bl{J}}{p, j \mapsto \upd{x[j] < x[p]}{j}\ \ \upd{x[p] \le x[j]}{p}}) \\
  &\quad\wedge\ \al{ks, Ns, Js}{k, n'', J \mapsto J = \cnt(k + 1, 1, n'')} \\
  &\quad\wedge\ \al{ks, Ns, js}{k, n'', j \mapsto j = \intf(n'') + k + 1}\\
  &\quad\wedge\ \al{ks, js}{k, j \mapsto n \le j} \\
  &\quad\wedge\ \al{ks, Js_2}{k, J \mapsto \al{\bl{J}}{j \mapsto j < n_2}}\ \wedge\ Js_2 = Js\ \wedge\ n_2 = n \\
  &\quad\Longrightarrow\ \cur(\tr{a}{ks, is, ms}{x, k, i, m \mapsto \swap{x}{k}{m}})\\
  &\qquad\qquad = \cur(\tr{a}{ks}{x, k \mapsto \swap{x}{k}{\argmin(x, k, n)}})
\end{align*}
The left-hand side is matched against $b = \cur(B)$ with the definition \eqref{eq:B} of $B$, which binds
$a$, $ks := \bl{K}$, $is := \bl{I}$, $ms := \seq{M}$. The side conditions are then searched among the
facts in order, each binding a new sequence: \eqref{eq:f1} binds $Ms := \seq{M}'$; \eqref{eq:f2} binds
$Js := \seq{J}'$, where the dependency $\bl{\tr{a}{\ldots}}$ of the pattern is matched against $\bl{B}$
through the definition of $B$; \eqref{eq:f4b} binds $Ns := \seq{N}$; \eqref{eq:f5b} binds $js := \seq{J}$;
\eqref{eq:f7} binds $n$. The remaining condition is the inner test \eqref{eq:f6}, which is matched by a
duplicate $Js_2$ with a fresh bound $n_2$, identified afterwards by the equations $Js_2 = Js$ and $n_2 = n$:
a side condition whose sequence variables are all bound is \emph{decided} by the simplifier rather than
matched against the facts, and the decision fails, since the facts keep $\bl{I}$ as a dependency and
rename the bound variables. The equations are closed and trivially decided.

Soundness of the lemma: given the side conditions, in iteration $k$ the selection trace $M$ scans the
positions $k + 1, \ldots, k + n''$ with $n'' + k + 1 \ge n$ (exit) and $k + 1 + n'' - 1 < n$ if $n'' > 0$
(test), hence exactly $k + 1, \ldots, n - 1$, so $\cur(M)$ is a position of the minimum of $x[k..n)$,
i.e. $\argmin(x, k, n)$ for any choice of $\argmin$ among the minima. The result is the trace of swaps of
$k$ with $\argmin(x, k, n)$, the \emph{summary} of the whole algorithm as a trace of an abstract step.
After the rewrite, $b = \cur(\tr{a}{\bl{K}}{x, k \mapsto \swap{x}{k}{\argmin(x, k, n)}})$.

Two workarounds are visible: the inner test is kept in its raw form \eqref{eq:f6} as a side condition
instead of being resolved to a bound on $n''$, since the resolved form $k + 1 + \intf(n'') \le n$ is
matched by \lemma{above\_iterations\_def} of \texttt{teq.dfy} (pattern $k < t$ with the dependency $t$
synthesized as a map), which hides the fact; and the initial value $k + 1$ and the final value
$\intf(n'') + k + 1$ are written literally in the normal form of the simplifier, since a higher-order
summand $X(k)$ matches a single summand only.

\subsection{Step 3: sortedness by induction, permutation by a generic lemma}\label{sec:sel:ind}

The postcondition $\sorted(b, 0, n)$ is now
\[ \sorted(\cur(\tr{a}{\bl{K}}{x, k \mapsto \swap{x}{k}{\argmin(x, k, n)}}), 0, n). \]
The invariant of
selection sort after $k$ iterations is that the prefix $x[0..k)$ is sorted and all of it is at most all
of $x[k..n)$, written with a second abstract predicate $\partitioned(x, lo, mid, hi)$ (all of $x[lo..mid)$
are at most all of $x[mid..hi)$). The induction is a lemma over a trace of arbitrary steps $F$:
\begin{align*}
  &\lemma{sorted\_prefix\_steps}:\\
  &\quad  n = \intf(m)\ \wedge\
    \big(\forall x, k.\ 0 \le k < n \wedge \sorted(x, 0, k) \wedge \partitioned(x, 0, k, n) \\
  &\qquad\qquad \Longrightarrow \sorted(F(x, k), 0, k + 1) \wedge \partitioned(F(x, k), 0, k + 1, n)\big) \\
  &\quad\Longrightarrow\ \sorted(\cur(\tr{a}{\bl{\cnt(0, 1, m)}}{x, k \mapsto F(x, k)}), 0, n)\ \Longleftrightarrow\ \mathsf{true}
\end{align*}
The left-hand side matches with $F := \lambda x\, k.\ \swap{x}{k}{\argmin(x, k, n)}$ and $m := n'$,
through the definition $K = \cnt(0, 1, n')$ in the dependency position. The condition $n = \intf(m)$
is the fact $\intf(n') = n$ of \eqref{eq:outer}; it fixes the number of steps to exactly $n$, since
further steps would be unconstrained by the second condition (a lemma with $n \le \intf(m)$ would be
unsound). The quantified condition is decided by the simplifier: $F$ is beta-reduced, and the two
conjuncts of the conclusion are rules on atoms,
\begin{align*}
  &\lemma{sorted\_swap\_argmin}:\\
  &\quad k < n \wedge \sorted(x, 0, k) \wedge \partitioned(x, 0, k, n)\\
    &\quad\Longrightarrow \sorted(\swap{x}{k}{\argmin(x, k, n)}, 0, k + 1) \Longleftrightarrow \mathsf{true},\\
  &\lemma{partitioned\_swap\_argmin}:\\
  &\quad k < n \wedge \sorted(x, 0, k) \wedge \partitioned(x, 0, k, n)\\
    &\quad\Longrightarrow \partitioned(\swap{x}{k}{\argmin(x, k, n)}, 0, k + 1, n) \Longleftrightarrow \mathsf{true},
\end{align*}
whose side conditions are the hypotheses of the quantified implication. The base case is inside the
induction lemma: $\sorted(a, 0, 0)$ and $\partitioned(a, 0, 0, n)$ hold trivially (two further lemmas,
\lemma{sorted\_empty} and \lemma{partitioned\_empty}, state them for completeness). This is the method of
the paper's \lemma{current\_trace\_bounded}, an invariant scheme as a lemma with a quantified side
condition: the lemma database proves the induction step, the shape of the invariant is the lemma's
knowledge of selection sort, and no invariant inference takes place.

The postcondition $\perm(b, a)$ follows from a generic lemma on traces of swaps at positions that may
depend on the state and the iteration (a higher-order pattern):
\begin{align*}
  &\lemma{perm\_swaps\_one}:\\
  &\quad \perm(\cur(\tr{a}{ks}{x, k \mapsto \swap{x}{P(x, k)}{Q(x, k)}}), c)\\
  &\quad\Longleftrightarrow \perm(a, c)
\end{align*}
with $P := \lambda x\, k.\ k$ and $Q := \lambda x\, k.\ \argmin(x, k, n)$, and then \lemma{perm\_refl}.
The whole goal closes with \texttt{./Cuvee.sh -teq -simplify}; a wrong postcondition, e.g.
$\sorted(b, 0, n + 1)$, leaves a goal.

\subsection{What is generic and what is specific}

Of the lemmas used, \lemma{invs\_sums\_counter}, \lemma{invs\_current\_counter} and
\lemma{perm\_swaps\_one} are independent of the algorithm; \lemma{sorted\_prefix\_steps} is the
induction for every algorithm that grows a sorted prefix (selection sort, and insertion sort in a
variant with two dependencies); the summary \lemma{current\_selection\_swaps} and the two step lemmas are the
algorithm's knowledge: what one inner loop computes, and why one step keeps the invariant. The abstract
predicates and $\argmin$ are axiomatized by these lemmas only; a definition in Lean, as for
\texttt{teq.dfy}, would make the step lemmas and the summary provable.
